\subsubsection{Konstrukcja zestawu testowego}
Protokół ewaluacji stosuje zadanie identyfikacji galeria-vs-sonda:

\paragraph{Zestaw galerii:} Skonstruowany z pierwszej połowy sesji (folderów) dla każdej osoby w podziale testowym. Dla każdej osoby osadzania z wielu czystych (nie okluzjowanych) obrazów są uśredniane:

\begin{equation}
\mathbf{e}_{\text{gallery}} = \frac{1}{n} \sum_{i=1}^{n} \text{L2-Normalize}(\mathbf{e}_i)
\end{equation}

Zapewnia to niezawodny prototyp osoby do dopasowania podobieństwa.

\paragraph{Zestaw sondy:} Skonstruowany z drugiej połowy sesji. Te obrazy są poddane syntetycznej okluzji podczas ewaluacji (szczegóły w sekcji \ref{subsubsec:eval_occlusion}).

\subsubsection{Ekstrakcja embeddingu i indeksowanie}
\begin{enumerate}
    \item Ekstrahuj embedding dla wszystkich obrazów galerii i sondy
    \item Normalizuj L2: $\hat{\mathbf{e}} = \frac{\mathbf{e}}{\|\mathbf{e}\|_2}$
    \item Indeksuj embedding galerii przy użyciu FAISS IndexFlatIP (Inner Product), gdzie podobieństwo cosinusowe na wektorach znormalizowanych L2 równa się iloczynowi wewnętrznemu
\end{enumerate}

\paragraph{Rozdział walidacji i ewaluacji:} w trakcie treningu stosowana jest pomocnicza walidacja na 500 parach (średnia cosinusowa czyste vs okluzjowane), próbkowanych z części testowej. Wyniki Rank-1/Rank-3 pochodzą z niezależnego pipeline’u FAISS dla zadania identyfikacji galeria-vs-sonda.

\subsubsection{Symulacja okluzji w ewaluacji}
\label{subsubsec:eval_occlusion}

Podczas ewaluacji każdy obraz sondy jest syntetycznie okluzjowany przy użyciu deterministycznego protokołu, aby zapewnić powtarzalność:

\begin{itemize}
    \item \textbf{Wysokość okluzji}: 20 pikseli
    \item \textbf{Szerokość okluzji}: jeśli dostępny \texttt{bbox} z adnotacji, pasek ograniczany jest do zakresu $[x_1, x_2]$; w przeciwnym razie pełna szerokość obrazu (0 do 112)
    \item \textbf{Kolor}: Czarny (0,0,0) [deterministyczny, w przeciwieństwie do losowej augmentacji treningowej]
    \item \textbf{Pozycja pionowa}: Region oczu, obliczony z punktów orientacyjnych:
    \begin{equation}
    c_y = \frac{y_{\text{lewe\_oko}} + y_{\text{prawe\_oko}}}{2}
    \end{equation}
    \item \textbf{Fallback}: Jeśli punkty orientacyjne niedostępne, $c_y = \frac{h}{2} - 10 = 46$ (dla 112×112)
    \item \textbf{Zakres Y}: $[c_y - 10, c_y + 10]$
\end{itemize}

Deterministyczna czarna okluzja (w przeciwieństwie do losowych kolorów w treningu) zapewnia rygorystyczną, powtarzalną ewaluację podczas testowania uogólnienia na wariacje kolorów.

\paragraph{Wyrównanie w ewaluacji:} obrazy testowe w zbiorze \texttt{webface\_112x112} są już wstępnie wyrównane.

\subsubsection{Wyszukiwanie podobieństwa i metryki}
Dla każdego obrazu sondy $p$:
\begin{enumerate}
    \item Zapytaj indeks FAISS z $k=3$ najbliższymi sąsiadami
    \item Pobierz top-3 tożsamości galerii ze wynikami podobieństwa: $(s_1, s_2, s_3)$
\end{enumerate}

\paragraph{Dokładność Rank-k:} 
\begin{equation}
\text{Rank-}k = \frac{\text{\# zapytań z prawidłową tożsamością w top-}k}{N_{\text{ogółem}} }
\end{equation}

W szczególności:
\begin{itemize}
    \item \textbf{Dokładność Rank-1}: $P(\text{tożsamość}[0] = \text{rzeczywistość})$
    \item \textbf{Dokładność Rank-3}: $P(\text{rzeczywistość} \in \{\text{tożsamość}[0,1,2]\})$
\end{itemize}

\subsubsection{Wyjście i logowanie}
Szczegółowe wyniki są logowane w formacie CSV z kolumnami:
\begin{lstlisting}
query_id, top1_id, sim1, top2_id, sim2, top3_id, sim3, found_in_top3
\end{lstlisting}